\chapter{نتیجه‌گیری و کار‌های آینده}

در این مقاله، یک رویکرد نوین برای ترکیب سیستم‌های چندعاملی، معماری ابری-مه-لبه و فناوری دوقلوی مجازی در زمینه مدیریت هوشمند فرآیندها ارائه شده است.
استفاده از عامل‌ها و مفاهیم داده‌های مرتبط معنایی در این رویکرد، به‌وضوح در پیاده‌سازی دوقلوی مجازی برای نه‌تنها صنایع تولید، بلکه برای کاربرد‌های دیگر نیز توصیف می‌شود. یکی از مزایای اصلی این رویکرد، توانایی تطبیق و تکامل خدمات تصمیم‌گیری دوقلوی مجازی با گذر زمان است که با افزودن جریان‌های داده جدید و به‌کارگیری تکنیک‌های مدل‌سازی نوین قابل انجام و پیشرفت است.

سبک معماری سیستم‌های چند عاملی در موفقیت این رویکرد نقش کلیدی داشته است زیرا امکان ترکیت بی‌وقفه عامل‌ها و میکروسرویس‌ها را فراهم می‌آورد. در این معماری، تلاش بر این است تا با همکاری بهتر عامل‌ها و تبادل اطلاعات بین آن‌ها، موانع موجود برای ایجاد دوقلوی مجازی از بین برود و نتیجه هر چه بهتری را به همراه داشته باشد.

آینده پژوهش می‌تواند به سمت ترکیب الگوریتم‌های هوشمند‌سازی مدرن حرکت کند. استفاده از روش‌های یادگیری توزیع‌شده مانند یادگیری مشارکتی\LTRfootnote{Federated Learning}
می‌تواند توانایی عامل‌ها را افزایش دهد. همچنین بررسی روش‌های مختلف ارتباط عامل‌ها و لایه‌های موجود در دوقلوی مجازی، می‌تواند مسیر مناسبی در گسترش این پژوهش باشد. علاوه‌براین، جزئیات هر عامل و لایه‌ها قابلیت شخصی‌سازی بر اساس نوع کاربرد معماری دارد و می‌توان در آن نو‌آوری‌های زیادی داشت.